\documentclass[12pt,a4paper]{article}
\usepackage[utf8]{inputenc}
\usepackage[indonesian]{babel}
\usepackage{amsmath}
\usepackage{amsfonts}
\usepackage{amssymb}
\usepackage{geometry}

\geometry{margin=2.5cm}

\title{Menghitung Sentroid dengan Teorema Pappus}
\author{Ahmad Fakhrul Bawani}
\date{}

\begin{document}

\maketitle

\section*{Soal}

Use Pappus's theorem for surface area and the fact that the surface area of a sphere of radius $\rho$ is $4\pi\rho^2$ to find the centroid of the semicircle $y = \sqrt{\rho^2 - x^2}$.

\subsection*{1. Teorema Pappus untuk Luas Permukaan}
Teorema Pappus untuk luas permukaan menyatakan bahwa jika sebuah kurva bidang datar dengan panjang busur $L$ diputar mengelilingi sebuah poros yang tidak memotong kurva tersebut, maka luas permukaan benda putar ($S$) yang dihasilkan adalah:
\begin{equation*}
  S = 2\pi \cdot \bar{y} \cdot L
\end{equation责任}
Di mana $\bar{y}$ adalah jarak tegak lurus dari pusat massa (sentroid) kurva ke kapak/poros putaran (dalam hal ini, sumbu-$x$).

\subsection*{2. Menentukan Karakteristik Kurva Setengah Lingkaran}
Kurva yang diberikan adalah $y = \sqrt{\rho^2 - x^2}$, yang merupakan busur setengah lingkaran bagian atas dengan jari-jari $\rho$.

\begin{itemize}
  \item \textbf{Simetri Koordinat Sentroid ($\bar{x}$):}
    Karena kurva setengah lingkaran $y = \sqrt{\rho^2 - x^2}$ simetris secara reflektif terhadap sumbu-$y$, maka posisi horizontal sentroidnya berada tepat di tengah, yaitu:
    \begin{equation*}
      \bar{x} = 0
    \end{equation*}
  \item \textbf{Panjang Busur Kurva ($L$):}
    Keliling lingkaran penuh berjejari $\rho$ adalah $2\pi\rho$. Karena kurva ini berbentuk setengah lingkaran, maka panjang busurnya adalah:
    \begin{equation*}
      L = \frac{1}{2} (2\pi\rho) = \pi\rho
    \end{equation*}
\end{itemize}

\subsection*{3. Aplikasi Perputaran dan Pencarian $\bar{y}$}
Jika kurva setengah lingkaran ini diputar $360^\circ$ mengelilingi sumbu-$x$, benda putar yang dihasilkan berbentuk sebuah **bola penuh** dengan jari-jari $\rho$.

Berdasarkan fakta geometri yang diberikan pada soal, luas permukaan bola tersebut adalah:
\begin{equation*}
  S = 4\pi\rho^2
\end{equation*}

Substitusikan nilai luas permukaan $S = 4\pi\rho^2$ dan panjang busur $L = \pi\rho$ ke dalam rumus Teorema Pappus:
\begin{align*}
  S &= 2\pi \cdot \bar{y} \cdot L \\
  4\pi\rho^2 &= 2\pi \cdot \bar{y} \cdot (\pi\rho) \\
  4\pi\rho^2 &= 2\pi^2\rho \cdot \bar{y}
\end{align*}

Selesaikan persamaan di atas untuk mencari nilai koordinat vertikal sentroid ($\bar{y}$):
\begin{align*}
  \bar{y} &= \frac{4\pi\rho^2}{2\pi^2\rho} \\
  \bar{y} &= \frac{2\rho}{\pi}
\end{align*}

\subsection*{Kesimpulan}
Jadi, dengan memanfaatkan Teorema Pappus untuk luas permukaan, sentroid dari kurva setengah lingkaran tersebut berada pada koordinat \textbf{$\left(0, \dfrac{2\rho}{\pi}\right)$}.

\end{document}